\documentclass[leqno, a4paper, 12pt]{jsarticle}
\usepackage{amssymb}
\usepackage{amsthm}
\usepackage{amsmath}
\usepackage{comment}
\usepackage{url}

\usepackage{enumitem}
\usepackage{showkeys}


%定理環境 thmはあまり使わずpropをなるべく使う。
%主張は基本的に証明の中で使い、そのため番号を付けない。
%注意環境を付けてみた。
\newtheorem{thm}{定理}
\newtheorem{prop}[thm]{命題}
\newtheorem{cor}[thm]{系}
\newtheorem{lem}[thm]{補題}
\newtheorem{df}[thm]{定義}
\newtheorem{exam}[thm]{例}
\newtheorem{fact}[thm]{事実}
\newtheorem*{claim}{主張}
\newtheorem*{cau}{注意}
\newtheorem*{rmk}{注意}
\renewcommand{\proofname}{証明}
%矢印たち。
%\toは\rightarrowと同じ。\getsは\leftarrowと同じ。同値は\iff
\newcommand{\To}{\Longrightarrow}
\newcommand{\Gets}{\LongLeftarrow}
%黒板太字
\newcommand{\nn}{\mathbb{N}}
\newcommand{\zz}{\mathbb{Z}}
\newcommand{\qq}{\mathbb{Q}}
\newcommand{\rr}{\mathbb{R}}
\newcommand{\cc}{\mathbb{C}}
%無限大
\newcommand{\mgn}{\infty}
%差集合は\setminusで統一する。等号の否定は\neq
%真部分集合は\subsetneqq　偏微分は\partial
%ディスプレイスタイル。\dfracでディスプレイ分数
\newcommand{\dpst}{\displaystyle}

\title{論文メモ
\large{\\--- エルデシュ的な話：代数？ ---}
}
\author{ナンブ　キトラ}
\date{\today}
\begin{document}
\maketitle

本棚を整理したいので，
溜まっている論文のメモを書いて未来への手紙とすることにする．

以下はなんかエルデシュっぽい数学で和に関する
話題を集めた論文の束だと思う．

すごく当たり前の話だが，
実数$\rr$を$\qq$上の代数
と考えた時の超越次数は連続体濃度になる．
ゆえに
選択公理を
仮定すれば連続体濃度をもつ実数の部分集合であって，
$\qq$上代数的独立なもの
をとることができる．
フォン・ノイマンの
論文
\cite{MR1512442}
では
超越基底ではないが，
とにかく連続体濃度をもつ
代数的独立な次数の部分集合を構成している．
正の実数$t$に関して
$A_{t}$
を
\[
A_{t}=\sum_{n=0}^{\infty}2^{2^{\lfloor tn \rfloor -2^{n^{2}}}}
\]
と定義すると
集合
$\{\, A_{t}\mid t\in (0, \infty)\, \}$
は代数的独立集合になるし，
$t$が異なると$A_{t}$も異なる．
これはもちろんZFの定理である．

論文
\cite{MR0983806}
は上のフォンノイマンの論文を読んで
自分でも代数的独立な実数を作りたいと思って
参考にした論文だと思う．


エルデシュとストーンの定理
\cite{MR0260958}
は
実数のボレル部分集合$A$
と$B$であって
$A+B=\{\, a+b\mid a\in A, b\in B\, \}$
が
ボレルにならない例を構成している．
特に彼らの集合はコンパクト集合と
$G_{\delta}$集合である．
コンパクト集合と閉集合の和が閉集合になることは
一般の位相群でも証明できる事柄である．
このことから特に一般の位相群でもコンパクト集合と
$F_{\sigma}$集合の和が$F_{\sigma}$
ということがわかる．
特に距離化可能な位相群であれば
コンパクト集合と開集合の和が
$F_{\sigma}$になることがわかる．
ということで彼らの構成した和がボレルに
ならないコンパクト集合と
$G_{\delta}$集合というのはわかっていることと
わかっていなかたことのちょうど境目である．

エルデシュとキューネンの論文
\cite{MR0641304}
は
ボレルではなく，もっと
測度論的話である．
特に$\rr$の二つの部分$\qq$線型空間
$G_{1}$, 
$G_{2}$であって，
$G_{1}+G_{2}=\rr$
でなおかつこれら二つの集合の測度
が$0$になるものを構成している．
ここら辺から集合論っぽい感じになる．

論文
\cite{MR1149411}
は上のエルデシュ・キューネンの論文を受けて，
何やらユークリッド空間のパラドキシカルな
分解について研究しているようだ．
バナッハ・タルスキーの話の仲間．


論文
\cite{MR1142923}
は普遍零集合の話らしい．
$\rr^{n}$の部分集合が普遍零集合
(universal null set)
であるとは
$\rr^{n}$上のルベーグ測度に関して連続
な任意のボレル測度$\mu$について零集合になることである．
よくわかんないけど，普遍零集合の和などを考察しているようだ．$Q$-set とか出てくる．


論文
\cite{MR2010332}
ではボレル集合と和の話を深掘りしているようだ．
$\sigma$-ideal
とか射影集合とか出てくる．

プレプリント
\cite{kuba2020cantor}
は
(一般の)カントール集合を含むような
$\rr$の部分体やベルンシュタイン集合を含むような部分体を
考察しているようだ．
%READ!
%myplain is the same to amsplain style. 
\nocite{*}
\bibliographystyle{myplaindoidoi}
\bibliography{eralg.bib}



\end{document}